%% kapa: exam=1 round=29 year=2018 period=1 form=A subject=civil_law subject_ko=민법 seq=8 number=8
%% source: Q-Net 감정평가사 기출문제 (cst00309) artlSeq=5202808 file=extracted/1008/5202808/2018년 제29회 감정평가사 1차시험 기출문제/2018년 제29회 감정평가사 1차시험 기출문제/1교시 A형.pdf
%% conversion: kordoc 4.24.51 -> tools/gen1.py (auto); review=proofread-sonnet; flags=none
\begin{question}{8}
불공정한 법률행위에 관한 설명으로 옳지 않은 것은? (다툼이 있으면 판례에 따름)
\begin{choices}
  \item 성립요건인 궁박, 경솔, 무경험은 그 중 하나만 갖추어도 충분하다.
  \item 궁박은 경제적 원인 외에 정신적 또는 심리적 원인에 기인할 수도 있다.
  \item 대리인에 의하여 행해진 법률행위에서 불공정한 법률행위가 문제되는 경우, 경솔이나 무경험은 대리인을 기준으로 판단한다.
  \item 무경험은 일반적인 생활체험의 부족을 의미하는 것으로, 어느 특정영역이 아니라 거래 일반에 대한 경험부족을 말한다.
  \item 매매계약이 약정된 매매대금의 과다로 인하여 불공정한 법률행위에 해당하는 경우, 무효행위의 전환에 관한 민법 제138조가 적용될 수 없다.
\end{choices}
\end{question}
\answer{5}
